An assumption is a belief of what you assume to be true in the future. You make assumptions based on your knowledge, experience or the information available on hand. These are anticipated events or circumstances that are expected to occur during your project's life cycle.

Assumptions are supposed to be true but do not necessarily end up being true. Sometimes they may turn out to be false, which can affect your project significantly. They add risks to the project because they may or may not be true. For example, if you are working on an outdoor unmanned vehicle, are you assuming that testing space will be available when needed? Are you relying on an external team or contractor to provide a certain subsystem on time? If you are working at a customer facility or deploying on their computing infrastructure, are you assuming you will be granted physical access or network credentials?

This section should contain a list of at least 5 of the most critical assumptions related to your project. For example:

The following list contains critical assumptions related to the implementation and testing of the project.

\begin{itemize}
    \item Steven McDermott will be available to review progress each sprint cycle
    \item Steven McDermott will provide a lab inventory list with item names and locations by the 2nd sprint cycle
    \item The university will provide a server to host the website and database by the 3rd sprint cycle
    \item The website will be permitted to sign users in with their UTA accounts
    \item A 3D scanner capable of capturing a full room will be available through purchase or university providing it
    \item The team will be granted access to scan Lab 208 and other UTA labs
    \item Wall space, power, and a network connection will be available for the kiosk inside Lab 208
\end{itemize}
